% concatenated from arxiv_18_06_2026.tex
\documentclass[12pt]{amsart}
\vfuzz2pt

\setlength{\arraycolsep}{0.17em}
\usepackage{amsmath}
\usepackage{amsfonts}
\usepackage{amssymb}
\usepackage{amsthm}
\usepackage{enumitem}
\usepackage{todonotes}
\usepackage{hyperref}

\newtheorem{df}{Definition}
\newtheorem{conj}{Conjecture}

\newtheorem{theorem}{Theorem}
\newtheorem{proposition}[theorem]{Proposition}
\newtheorem{lemma}[theorem]{Lemma}
\newtheorem{kor}[theorem]{Corollary}
\theoremstyle{definition}
\newtheorem*{example}{Example}
\newtheorem*{rem}{Remark}

\begin{document}
\title[On Diophantine pairs and triples of triangular numbers]{On Diophantine pairs and triples of triangular numbers}
\author{Marija {Bliznac Trebje\v{s}anin}}
\address{Faculty of Science, University of Split, Ru\dj{}era Bo\v{s}kovi\'{c}a 33, 21000 Split, Croatia}
\email{marbli@pmfst.hr}

\begin{abstract}
We investigate Diophantine pairs and triples of triangular numbers with the property $D(a)$ for a non-zero integer $a$. We prove that if a triangular number belongs to a $D(a)$-pair, it can be extended to infinitely many $D(a)$-triples of triangular numbers. Additionally, we determine infinite families of integers $a$ that admit such pairs, as well as families for which no $D(a)$-pairs can exist.
\end{abstract}

\maketitle

\noindent 2020 {\it Mathematics Subject Classification:}  11D09, 11B37
\\ \noindent Keywords: triangular numbers, recurrences, Diophantine m-tuples, generalized Pellian equations

\section{Introduction}

Let $a\neq0$ be an integer. We call a set of $m$ distinct positive integers a $D(a)$-$m$-tuple (or a Diophantine $m$-tuple with the property $D(a)$) if the product of any two of its distinct elements increased by $a$ is a perfect square. 
When $a=1$, the classical case introduced by Diophantus, such a set is simply referred to as a Diophantine $m$-tuple.

The study of such sets traces back to Fermat, who discovered the first Diophantine quadruple $\{1, 3, 8, 120\}$. A central question in this area of research is determining the maximum possible size of these sets. It is now established that $D(1)$-quintuples \cite{htz}, $D(4)$-quintuples \cite{petorke_btf}, and $D(-1)$-quadruples \cite{bcm} do not exist. Another prominent direction of research involves restricting the elements of $D(a)$-$m$-tuples to specific integer sequences, such as Fibonacci numbers (see, for instance, \cite{duje_gen, fujita_luca_fib}). For a comprehensive overview of various generalizations and related problems, we refer the reader to \cite[Section 1.5]{duje_book}.

For a positive integer $n$, the $n$-th triangular number, denoted by $T_n$, is defined as the sum of the first $n$ positive integers, given by the explicit formula
$$T_n=\frac{n(n+1)}{2}.$$

From the main result of \cite{mnde}, it follows that  
$\{T_n,T_{n+4},T_{4n^2+20n+8}\}$  
is a Diophantine triple of triangular numbers. In \cite{mbt_triangular}, the author presented a construction showing that for any positive integers $k$ and $n$, the triangular number $T_n$ is a member of infinitely many $D(k^2)$-triples containing only triangular numbers.

Let $a\neq 0$ and $n\geq 1$ be integers. We are looking for positive integers $m$ and $r$ satisfying
$$T_nT_m+a=r^2.$$
This equation can be rewritten as a generalized Pellian equation 
\begin{equation}\label{pellova_x}
    x^2-n(n+1)y^2=16a-n(n+1),
\end{equation}
where the unknowns are $x=4r$ and $y=2m+1$. An analysis of the properties of this equation will provide the main results concerning the existence of $D(a)$-pairs and $D(a)$-triples of triangular numbers. 

The existence of a solution to equation \eqref{pellova_x}, with an odd integer $y$, depends on properties of both $a$ and $n$. As shown in \cite{mbt_triangular}, for $a=k^2$, $k\in\mathbb{N}$, we have a solution for every integer $n\geq 1$. Determining the conditions on $a$ and $n$ under which a solution exists will be the primary focus of Sections \ref{section:ex_pairs} and \ref{section:member_pairs}. In the general case for an integer $a\neq 0$, by analyzing the sequences $(m_k)_{k \in \mathbb{N}}$ arising from the solutions to equation \eqref{pellova_x}, we will prove that for any two successive terms of the sequence, the set $\{T_{m_k}, T_{m_{k+1}}\}$ forms a $D(a)$-pair. This naturally leads to our main result.

\begin{theorem}\label{tm_2}
Let $a\neq 0$ and $n\geq 1$ be integers such that there exists a positive integer $m$ for which $\{T_n,T_m\}$ is a $D(a)$-pair. Then $T_n$ is an element of infinitely many $D(a)$-triples of triangular numbers.
\end{theorem}

The remainder of this paper is organized as follows. In Section \ref{section:triple}, we cover the basic properties of the associated generalized Pellian equation and prove Theorem \ref{tm_2}. In Section \ref{section:ex_pairs}, we explore the conditions on the integer $a$ for which $D(a)$-pairs can or cannot exist, providing infinite families for both cases. Finally, in Section \ref{section:member_pairs}, we investigate the necessary conditions an integer $n$ must satisfy for $T_n$ to belong to a $D(a)$-pair, with a specific focus on the case $a=-1$, and describe an algorithm to generate all such values up to a given bound.

\section{Generalized Pellian equation and existence of $D(a)$-triples}\label{section:triple}

Let $a\neq 0$ and $n\geq 1$ be integers. We are searching for positive integers $m$ and $r$ such that $T_n \cdot T_m +a= r^2$, or more precisely,
$$\frac{n(n+1)}{2} \cdot \frac{m(m+1)}{2}+a = r^2.$$
After some algebraic manipulation, this yields the generalized Pellian equation
\begin{equation}\label{pellova_m}
    (4r)^2 - n(n+1)(2m+1)^2 = 16a-n(n+1),
\end{equation}
which can be written in the simpler form of the equation \eqref{pellova_x} in terms of the variables $x$ and $y$. 

From the theory of generalized Pellian equations (see \cite[Theorem 109]{nagell}), we know that if equation \eqref{pellova_x} has a solution, it has a finite number of classes of solutions. As described in \cite[Theorem 108a]{nagell}, each class is generated by a fundamental solution $(x^*,y^*)$ satisfying the inequalities
\begin{equation}\label{ineq_granice_fund}
0<y^*\leq \frac{Y_1}{\sqrt{2(X_1-1)}}\sqrt{|N|},\quad 0\leq|x^*|\leq\sqrt{\frac{1}{2}(X_1-1)\cdot |N|},
\end{equation}
where $N=16a-n(n+1)$ and $(X_1,Y_1)=(2n+1,2)$ is the fundamental solution of the corresponding Pell equation 
\begin{equation}\label{eq_pellova_1}
X^2-n(n+1)Y^2=1.
\end{equation}

We consider a specific class of solutions to equation \eqref{pellova_x} and denote its fundamental solution by $(x^*,y^*)$. Then every positive integer solution $(x,y)$ belonging to this class can be expressed as
$$x+y\sqrt{n(n+1)}=(x^*+y^*\sqrt{n(n+1)})\left(X_1+Y_1\sqrt{n(n+1)}\right)^k,$$
for some integer $k \geq 0$ if $x^* > 0$, and $k \geq 1$ if $x^* \leq 0$. 

For generalized Pellian equation  \eqref{pellova_x}, the following relation holds: 
\begin{equation}\label{niz_y_x}
    y_{k+1}=X_1y_k+Y_1x_k=(2n+1)y_k+2x_k,\quad k\geq 0.
\end{equation}


\begin{lemma}\label{lema_svojstva_xy}
    If $(x,y)$ is a solution to the equation \eqref{pellova_x} with $y$ an odd number, then  $x$ is divisible by $4$ and $y^*\geq 1$ is also an odd number. Moreover, if $a$ is not a perfect square, then $y^*\geq 3$. 
\end{lemma}

\begin{proof} 
Equation \eqref{pellova_x} can be rewritten as $x^2 - 16a = n(n+1)(y^2 - 1)$. 

Since $y$ is odd, $y^2 - 1$ is a multiple of $8$ and $n(n+1)(y^2 - 1)$ is divisible by $16$, implying that $x$ is divisible by $4$.

Finally,  the assumption $y=1$ leads to $x^2 = 16a$, implying that $a$ is a perfect square. 
\end{proof}

For $a=-1$ it holds that equations \eqref{pellova_m} and \eqref{pellova_x} are equivalent, meaning that the positive integer solutions $(x,y)$ of equation \eqref{pellova_x} are always of the form $(4r,2m+1)$ for some positive integers $r$ and $m$.  

\begin{kor}\label{kor_svojstva_xy_minus_1}
    Let $a=-1$. If $(x,y)$ is a solution to the equation \eqref{pellova_x} then $y$ is an odd number and $x$ is divisible by $4$. Moreover, if $y$ is a positive integer, then $y\geq 3$. 
\end{kor}


The sequence of solutions for the second coordinate, $(y_k)_{k \geq 0}$, forms a linear recurrence sequence 
\begin{align}\label{eq_rek_y}
&y_0=y^*,\quad y_1=2x^*+(2n+1)y^*,\nonumber\\
&y_{k}=2(2n+1)y_{k-1}-y_{k-2},\ k \geq 2.
\end{align}
Since $y_k=2m_k+1$, we can express the sequence of indices $(m_k)$ such that $\{T_n,T_{m_k}\}$ is a $D(a)$-pair as follows:
\begin{align}\label{eq_niz_mk}
&m_0=\frac{1}{2}(y^*-1),\quad m_1=\frac{1}{2}(2x^*+(2n+1)y^*-1),\nonumber\\
 &m_{k}=2(2n+1)m_{k-1}-m_{k-2}+2n,\quad k\geq 2.
\end{align}
Note that if $x^* =0$, the value $r=2|x^*|$ is not a positive integer, hence, we should only consider the indices $m_k$ for $k \geq 1$.

\begin{lemma}\label{lem_y}
    For the sequence $(y_k)_{k \geq 0}$ of solutions to the generalized Pellian equation \eqref{pellova_x}, the following identity holds:
$$y_{k+1}^2 - 2(2n+1)y_k y_{k+1} + y_k^2 =  64a - 4n(n+1).$$
\end{lemma}
\begin{proof}
From equation \eqref{niz_y_x}, we can express $2x_k$ as $2x_k = y_{k+1} - (2n+1)y_k$. Squaring this identity and applying the relation $x_k^2=n(n+1)(y_k^2-1)+16a$ yields the desired equality. 
\end{proof}

Now we are ready to prove that $\{T_{m_k},T_{m_{k+1}}\}$ forms a $D(a)$-pair for all $k\geq 1$. 

\begin{proposition}
    Let $(m_k)_{k\geq 1}$ be the sequence defined in \eqref{eq_niz_mk}. Then 
    $$T_{m_k}T_{m_{k+1}}+a$$
    is a square of a positive integer. 
\end{proposition}
\begin{proof}
    Notice that $8T_n+1=(2n+1)^2$ holds for every integer $n$. Then 
    \begin{align*}
        T_{m_k}T_{m_{k+1}}+a&=\frac{(2m_k+1)^2-1}{8}\cdot\frac{(2m_{k+1}+1)^2-1}{8} +a\\
        &=\frac{y_k^2-1}{8}\cdot\frac{y_{k+1}^2-1}{8} +a=\frac{y_k^2y_{k+1}^2-y_k^2-y_{k+1}^2+1+64a}{64}.
    \end{align*}
    By substituting the identity from Lemma \ref{lem_y}, we obtain
    \begin{align*}
        T_{m_k}T_{m_{k+1}}+a
        &=\frac{y_k^2y_{k+1}^2+4n(n+1)+64-2(2n+1)y_ky_{k+1}+1+64a}{64}\\
        &=\left(\frac{y_ky_{k+1}-(2n+1)}{8}\right)^2\\
        &=\left(\frac{2m_km_{k+1}+m_k+m_{k+1}-n}{4}\right)^2.
    \end{align*}
    It remains to show that the expression inside the parentheses is an integer. It is easy to prove by mathematical induction that $y_ky_{k+1}-(2n+1)$ is divisible by $8$ for each $k\geq 0$. Properties from Lemma \ref{lema_svojstva_xy} are used in the proof. 
\end{proof}

This implies that any $T_n$ for which the associated equation \eqref{pellova_m} has a solution is a member of infinitely many $D(a)$-triples $\{T_n,T_{m_k},T_{m_{k+1}}\}$ for $k\geq 1$. Thus, we conclude the proof of Theorem \ref{tm_2}.


\section{Existence of $D(a)$-pairs}\label{section:ex_pairs}

In this section, we will explore for which positive integers $a\neq 0$ the equation \eqref{pellova_x} can have a solution. 

\begin{theorem}\label{thm:modular_9}
Let $a\neq 0$ be a positive integer such that there exists a $D(a)$-pair of triangular numbers. Then $a \not\equiv 2, 5 \pmod 9$.
\end{theorem}
\begin{proof}
    Let $r$ be a positive integer such that
    $$T_n\cdot T_m+a=r^2.$$
    Then $r^2\equiv 0,1,4,7\pmod 9$ and $T_n\cdot T_m\equiv 0,1,3,6\pmod 9$. From $a\equiv r^2- T_n\cdot T_m$, by considering all possibilities, we get 
    $a\equiv 0,1,3,4,6,7,8\pmod 9.$
\end{proof}

We will now prove that the converse does not hold. For each remaining class modulo $9$, we give an example of an integer $a$ that admits no $D(a)$-pair of triangular numbers. For illustrative purposes, we detail the proofs for $a=12$ and $a=-37$ ($a\equiv 3,8\pmod 9$). In the cases $a\in\{252,487,-158,42,52\}$, i.e., $a\equiv 0,1,4,6,7\pmod 9$, the proofs proceed analogously and are omitted for brevity. 

\begin{proposition} \label{prop:nonexist}
Let $a \in \{12, -37\}$. For any positive integer $n$, the equation \eqref{pellova_x}
does not have any integer solutions $(x, y)$ such that $y$ is odd. Consequently, there exists no pair of triangular numbers $\{T_n, T_m\}$ with the property $D(12)$ or $D(-37)$.
\end{proposition}

\begin{proof}
Assume, for the sake of contradiction, that there exists a positive integer $n$ such that equation \eqref{pellova_x} has an integer solution $(x, y)$ with an odd integer $y \ge 1$. 
By setting $y = 2m+1$ for some positive integer $m$, equation \eqref{pellova_x} can be rewritten as:
$$ x^2 = 16a + n(n+1)((2m+1)^2-1). $$
Notice that $n(n+1) = \frac{N_0^2-1}{4}$. Hence, it is useful to introduce the notation $N_0 = 2n+1$, $M_0 = 2m+1$ and $c = 2|x|$, multiply the equation by $4$ and obtain the following symmetric formulation:
\begin{equation} \label{eq:sym_combined_mixed}
c^2 = (N_0^2-1)(M_0^2-1) + 64a.
\end{equation}
By symmetry, we may assume without loss of generality that $m \geq n$, which implies $M_0 \geq N_0$. In what follows, we restrict our analysis of equation \eqref{eq:sym_combined_mixed} to this case.

If $M_0 = 1$ (or $N_0 = 1$), then  $c^2 = 64a$. 
 For $a=12$, this yields $c^2 = 768$, which is not a perfect square.
    For $a=-37$, this yields $c^2 = -2368$, which is trivially impossible.

Hence, any valid solution to \eqref{eq:sym_combined_mixed} requires $N_0 \ge 3$ and $M_0 \ge 3$. 

Equation \eqref{eq:sym_combined_mixed} can be rewritten as a generalized Pellian equation in variables $X=c$ and $Y=M_0$:
$$ X^2 - (N_0^2-1)Y^2 = 64a - (N_0^2-1). $$
If it admits any solutions with $M_0 \ge N_0 \ge 3$, then by the well-ordering principle, there must exist a solution for a fixed $N_0$ such that the positive integer $M_0$ is minimal. 
Let $(c, M_0)$ be this minimal integer solution. The fundamental solution $(N_0, 1)$ to the associated Pell equation $X^2 - (N_0^2-1)Y^2 = 1$ allows us to define a new solution $(c', M_0')$:
$$ (c', M_0') := \bigl(N_0 c - (N_0^2-1)M_0,\;\; N_0 M_0 - c\bigr). $$
Observe that the parity of the new elements $(c', M_0')$ is preserved, and that $(c', |M_0'|)$ is also a solution to equation \eqref{pellova_x}. 
Furthermore, $|M_0'| \neq 1$; otherwise, we would have either $16a=(m-n)^2$ or $16a=(m+n+1)^2$, which occurs only when $a$ is a perfect square. Since $M_0'$ is odd, we conclude that $|M_0'| \ge 3$.

For the new solution $(c', |M_0'|)$ to satisfy the  inequality $|M_0'| < M_0$, the value of $c$ must fall strictly between $(N_0-1)M_0$ and $(N_0+1)M_0$. Squaring this condition yields two inequalities that must hold simultaneously:
\begin{align}
    c^2 &< (N_0+1)^2 M_0^2 \implies 2M_0^2(N_0+1) > 64a + 1 - N_0^2, \label{upper_bound} \\
    c^2 &> (N_0-1)^2 M_0^2 \implies 2M_0^2(N_0-1) > -64a - 1 + N_0^2. \label{lower_bound}
\end{align}
If both bounds hold, we obtain a new integer solution with $3 \le |M_0'| < M_0$, which contradicts our assumption that $M_0$ is minimal. Therefore,  at least one of these inequalities must fail. This restricts the minimal solution to a finite set of exceptional pairs $(N_0, M_0)$.


\textbf{Case $a=12$:} For a positive parameter, \eqref{lower_bound} holds trivially. Inequality \eqref{upper_bound} fails when $2M_0^2(N_0+1) \le 769 - N_0^2$. This yields exactly six pairs with $M_0 \ge N_0 \ge 3$: $(3,3)$, $(3,5)$, $(3,7)$, $(3,9)$, $(5,5)$, and $(5,7)$. Substituting these into \eqref{eq:sym_combined_mixed} generates values for $c^2$ which are not perfect squares.


 \textbf{Case $a=-37$:} For a negative parameter, inequality \eqref{upper_bound} holds trivially. The lower bound does not hold when $2M_0^2(N_0-1) \le 2367 + N_0^2$. From equation \eqref{eq:sym_combined_mixed}, the condition $c^2 \ge 0$ imply that $(N_0^2 - 1) (M_0^2 - 1) \geq 2368$. Combining these two inequalities yields 13 pairs: $(3,19)$, $(3,21)$, $(3,23)$, $(5,11)$, $(5,13)$, $(5,15)$, $(5,17)$, $(7,9)$, $(7,11)$, $(7,13)$, $(9,9)$, $(9,11)$, and $(11,11)$. However, none of these pairs yields a perfect square for $c^2$.

Since the supposed minimal solution $(c, M_0)$ must belong to this set of pairs, and none of these pairs yield a perfect square for $c^2$, such a minimal solution cannot exist. Consequently, the equation has no solutions for $a=12$ or $a=-37$.
\end{proof}

\begin{rem}
It is worth noting that applying the same algorithm for $a=-1$ leads to the pair $(N_0, M_0) = (3,3)$. Substituting these values into equation \eqref{eq:sym_combined_mixed} yields $c=0$. This pair corresponds to $n=1$ and $m=1$, which exactly produces the fundamental solution $(x^*, y^*) = (0,3)$ for the initial generalized Pellian equation \eqref{pellova_x}. Since $m=n=1$, this trivial case does not yield a $D(-1)$-pair of distinct triangular numbers. However, we can generate an index $m>n$ by applying the recurrence relation \eqref{eq_niz_mk}, and we obtain $m_1 = 4$, i.e., the $D(-1)$-pair $\{T_1, T_4\}$.
\end{rem}


While it is known from \cite{mbt_triangular} that $D(a)$-pairs of triangular numbers exist whenever $a$ is a perfect square, we now construct additional infinite families of integers $a$ that admit such pairs. 

Using the well-known property that the sum of two consecutive triangular numbers is a perfect square, $T_{k+1} + T_k = (k+1)^2$, $k\geq 1$, we can construct explicit $D(a)$-pairs by fixing $a$ in relation to the sequence of triangular numbers. Thus the next proposition gives us infinite families of both positive and negative numbers $a$ for which we know $D(a)$-pairs of triangular numbers exist.

\begin{proposition} \label{prop:a_is_triangular}
Let $k \ge 1$ be an integer.
\begin{enumerate}
    \item If $a = T_k$, then the set $\{T_1, T_{k+1}\}$ is a $D(a)$-pair of triangular numbers.
    \item If $a = 1 - T_k$, then the set $\{T_2, T_{k+1}\}$ is a $D(a)$-pair of triangular numbers.
\end{enumerate}
\end{proposition}
\begin{proof}
Let $k\geq 1$ and $a = T_k$. Since $T_1 = 1$, we have:
$$
T_1 \cdot T_{k+1} + a = 1 \cdot \frac{(k+1)(k+2)}{2} + \frac{k(k+1)}{2}= (k+1)^2. 
$$
Thus, $\{T_1, T_{k+1}\}$ forms a $D(T_k)$-pair.

Now, let $a = 1 - T_k$. Substituting $T_2 = 3$, we have:
$$ T_2 \cdot T_{k+1} + a = 3 \cdot \frac{(k+1)(k+2)}{2} + 1 - \frac{k(k+1)}{2}=(k+2)^2. $$
Thus, $\{T_2, T_{k+1}\}$ is a $D(1 - T_k)$-pair for all $k \ge 1$.
\end{proof}

\section{Finding triangular numbers which form a $D(a)$-pair}\label{section:member_pairs}

In \cite{mbt_triangular} it has been shown that for every integer $n\geq 1$ and $a=k^2$, $k\geq 1$, triangular number $T_n$ is a member of a $D(a)$-pair of triangular numbers. It is natural to explore this further and for a fixed $a\neq0$, ask which triangular numbers $T_n$ can be members of a $D(a)$-pair of triangular numbers.

\begin{theorem} \label{thm:generalized}
Let $a$ be a non-zero integer. The triangular number $T_n$ can be a member of a $D(a)$-pair $\{T_n,T_m\}$, $m\geq 1$, only if all prime factors of $n(n+1)$ belong to the following set: the prime $2$, prime divisors of $a$, and odd primes $p$ for which $a$ is a quadratic residue modulo $p$. 
\end{theorem}
\begin{proof}
Assume that $n$ is a positive integer such that there exists a $D(a)$-pair $\{T_n,T_m\}$. Then   integers $x=4r$ and $y=2m+1$ are solutions of the equation \eqref{pellova_x}. Let $p$ be an arbitrary odd prime divisor of $n(n+1)$. 

After reducing equation \eqref{pellova_x} modulo $p$, we have:
$$ x^2 \equiv 16a \pmod p. $$
If $p \mid a$, the congruence trivially becomes $x^2 \equiv 0 \pmod p$, which is satisfied for all $x$ such that $p \mid x$. 

If $p \nmid a$, then $16a$ is not divisible by $p$. We can divide the congruence by $16$, which gives:
$$ r^2=(x/ 4)^2 \equiv a \pmod p. $$
This implies that $a$ is a quadratic residue modulo $p$. 
\end{proof}

Let's illustrate this condition further on $a=-1$.

\begin{kor}\label{kor:minus1}
Let $n$ be a positive integer such that there exists a positive integer $m$ for which $\{T_n,T_m\}$ is a $D(-1)$-pair. Then both $n$ and $n+1$ are numbers of the form $2^e\cdot p_1^{\alpha_1}\cdots p_l^{\alpha_l}$, where $p_i$ are prime numbers and $p_i\equiv 1\pmod{4}$ holds for each $i\in\{1,\dots,l\}$, $\alpha_1,\dots,\alpha_l$ are positive integers and $e\in\{0,2\}$ for $n$ and $e\in\{0,1\}$ for $n+1$.
\end{kor}
\begin{proof}
Let $a=-1$. From Theorem \ref{thm:generalized} we know that any odd prime factor $p$ of either $n$ or $n+1$ must be such that $-1$ is a quadratic residue modulo $p$. Hence, $p$ must be congruent to $1 \pmod{4}$.

Consequently, neither $n$ nor $n+1$ can be congruent to $3 \pmod{4}$. Since $n$ and $n+1$ are consecutive integers, the only possible congruences for the pair $(n, n+1)$ modulo $4$ are $(1, 2)$ or $(0, 1)$. This immediately eliminates $e=1$ for $n$ (since $n \not\equiv 2 \pmod{4}$) and $e \geq 2$ for $n+1$ (since $n+1 \not\equiv 0 \pmod{4}$). Thus, $e \in \{0, 1\}$ for $n+1$, and $e=0$ or $e \geq 2$ for $n$.

It remains to show that $e \leq 2$ for $n$. Assume, for the sake of contradiction, that $e \geq 3$, which implies $n \equiv 0 \pmod{8}$. From the definition of a $D(-1)$-pair, we have
$$ T_n T_m = r^2 + 1. $$
If $n \equiv 0 \pmod{8}$, then $T_n=\frac{n(n+1)}{2}$ is divisible by $4$. This implies that $r^2 \equiv 3 \pmod{4}$, which is impossible for any integer $r$. Thus, $n$ cannot be a multiple of $8$, proving $e \in \{0, 2\}$ for $n$.
\end{proof}

\begin{rem} 
 Let us show that the converse of Corollary \ref{kor:minus1}, and consequently Theorem \ref{thm:generalized}, does not hold. Number $n=100$ satisfies the necessary conditions of Corollary \ref{kor:minus1}. But there does not exist a pair of integers $(x^*,y^*)$ satisfying the bounds \eqref{ineq_granice_fund}. Hence, the triangular number $T_{100}$ is not a member of any $D(-1)$-pair containing only triangular numbers. It can, however, form a $D(-1)$-pair where the second element is not a triangular number, e.g., $\{T_{100}, 5653\}$ is a $D(-1)$-pair.
\end{rem}

Observe that from inequalities \eqref{ineq_granice_fund} and the fact that $(X_1,Y_1)=(2n+1,2)$, it follows that $y^*<2n+1$ for $n\geq \sqrt[3]{4|a|}$. From this observation and Lemma \ref{lema_svojstva_xy} we see that the fundamental solution with odd integer $y^*\geq 3$ yields a positive integer $m$ strictly smaller than $n$ such that $\{T_m,T_n\}$ is a $D(a)$-pair. This proves the next result. 

  \begin{proposition}
    Let $a\neq 0$ be an integer that is not a perfect square. 
    If a triangular number $T_n$, with $n \geq \sqrt[3]{4|a|}$, is a member of a $D(a)$-pair of triangular numbers, then there exists an integer $n_0$, $1 \leq n_0 < n$, such that the pair $(x,2n+1)$ is a solution to the equation $$x^2 - n_0(n_0+1)y^2 = 16a - n_0(n_0+1).$$
\end{proposition}
\medskip 

In the case $a=-1$, this holds for $n\geq 2$. 
 We exploit this property to illustrate an algorithm that systematically identifies all triangular numbers $T_n$ belonging to a $D(-1)$-pair of triangular numbers for $n \leq C$. Starting from the initial value $n=1$, we generate all corresponding solutions $m \leq C$ to the equation \eqref{pellova_m}. By iterating this process and treating each new solution as a seed for subsequent steps, we obtain all such values until they exceed the upper bound $C$.


\begin{example}
Let us set the upper bound to $C=1000$ and search for indices $n$ of members of $D(-1)$-pairs of triangular numbers. Starting with the base case $n=1$, we recursively generate the solutions. 

\begin{table}[h!]
\centering
\caption{Generation of indices $m \leq 1000$ }
\label{table_iteracije}
\renewcommand{\arraystretch}{1.3}
\begin{tabular}{|c|c|l|}
\hline
\textbf{Step} & \textbf{Seed $n$} & \textbf{New solutions $m \in (n, 1000]$} \\
\hline
1 & $1$ & $4, 25, 148, 865$ \\
\hline
2 & $4$ & $25, 457$ \\
\hline
3 & $25$ & $148, 457$ \\
\hline
4 & $148$ & $865$ \\
\hline
5 & $457$ & None \\
\hline
6 & $865$ & None \\
\hline
\end{tabular}
\end{table}

The algorithm naturally terminates at this point, as applying the recurrence relations to $n=457$ and $n=865$ yields solutions that exceed the bound of $1000$. Consequently, the complete set of numbers $n$ up to $1000$ such that $T_n$ is a member of a $D(-1)$-pair of triangular numbers is $\{1, 4, 25, 148, 457, 865\}$. 
\end{example}


\section*{Acknowledgment}
{The author is supported by the Croatian Science Foundation grant no.~IP-2022-10-5008 and by the University of Split, grant no.~IP-UNIST-44, funded by the European Union -- NextGenerationEU.}

\section*{Declarations}

\textbf{Data Availability} \\
Data sharing is not applicable to this article as no datasets were generated or analysed during the current study.

\vspace{1em}

\textbf{Use of Artificial Intelligence} \\
%\section*{Declaration of AI usage}
Gemini 3.1 Pro was used to explore research directions and refine language. All core concepts and mathematical proofs are the author's independent work. Claude Opus 4.8 generated an initial outline for Proposition \ref{prop:nonexist}, which the author fully corrected, formalized, and verified due to logical flaws in the AI output. The author assumes complete responsibility for the final manuscript's correctness.

\begin{thebibliography}{99}

\bibitem{mbt_triangular} M.~Bliznac Trebje\v{s}anin, {\it On Diophantine triples containing a triangular number}, {Math.~Slovaca} (2026). \url{https://doi.org/10.1515/ms-2025-1077}

\bibitem{petorke_btf} M.~Bliznac Trebje\v{s}anin, A.~Filipin, {\it Nonexistence of $D(4)$-quintuples}, {  J.~Number Theory}, \textbf{194} (2019), 170--217.

\bibitem{bcm} N.~C.~Bonciocat, M.~Cipu and M.~Mignotte, {\it There is no Diophantine D(-1)-quadruple}, J. London Math. Soc. \textbf{105} (2022), 63--99.

\bibitem{mnde}M.~N.~Deshpande, {\it One property of triangular numbers}, Portugaliae Math. \textbf{55} (1998), 381--383.

\bibitem{duje_gen} A. Dujella, {\it Generalization of a problem of Diophantus}, Acta Arith. \textbf{65} (1993), 15--27.

\bibitem{duje_book} A. Dujella, {\it Diophantine $m$-tuples and Elliptic Curves}, Springer, Cham, 2024.

\bibitem{fujita_luca_fib} Y.~Fujita, F.~Luca, {\it There are no Diophantine quadruples of Fibonacci numbers}, Acta Arith. \textbf{185} (2018), 19--38.

\bibitem{htz} B.~He, A.~Togb\'{e}, V.~Ziegler, {\it There is no Diophantine quintuple}, {  Trans.~Amer.~Math.~Soc.~} \textbf{371} (2019), 6665--6709.

\bibitem{nagell}
T.~Nagell, \emph{Introduction to Number Theory}, 2nd ed., Chelsea Publishing Company, 1964.

\end{thebibliography}

\end{document}
